% SwartzNet whitepaper — LaTeX source.
% GENERATED from docs/whitepaper.md by docs/build-whitepaper.sh (prose is
% authoritative in the .md). Self-contained: compiles with a standard TeX Live
% via `latexmk -pdf whitepaper` or `make -C docs whitepaper`.
\documentclass[11pt]{article}

% ---- fonts & typography ---------------------------------------------------
\usepackage[T1]{fontenc}
\usepackage[utf8]{inputenc}
\usepackage{newtxtext}          % Times-like professional serif text
\usepackage{newtxmath}
\usepackage{microtype}          % refined justification & character protrusion

% ---- page geometry --------------------------------------------------------
\usepackage[letterpaper,margin=1in]{geometry}
\linespread{1.04}
\setlength{\parindent}{1.4em}

% ---- lists & section headings ---------------------------------------------
\usepackage{enumitem}
\setlist{topsep=4pt,itemsep=2pt,parsep=0pt}
\usepackage{titlesec}
\setcounter{secnumdepth}{0}     % headings carry their own manual numbers (1., 2., ...)
\titleformat{\section}{\normalfont\large\bfseries}{}{0pt}{}
\titlespacing*{\section}{0pt}{1.5em}{0.55em}

% ---- links (clean: no colored boxes, working PDF bookmarks) ---------------
\usepackage[hidelinks]{hyperref}
\usepackage{xurl}               % breakable file paths in the references (\path)
\hypersetup{
  pdftitle={SwartzNet: Full-Text Search over the Mainline BitTorrent Network},
  pdfauthor={The SwartzNet Project},
  pdfsubject={A mainline-compatible distributed full-text search layer for BitTorrent},
  pdfkeywords={BitTorrent, DHT, full-text search, BEP-44, peer-to-peer, set reconciliation, RIBLT}
}

% pandoc helper (tight lists)
\providecommand{\tightlist}{\setlength{\itemsep}{0pt}\setlength{\parskip}{0pt}}

\begin{document}

% ---- title block ----------------------------------------------------------
\begin{center}
  {\LARGE\bfseries SwartzNet\par}
  \vspace{0.35em}
  {\large\bfseries Full-Text Search over the Mainline BitTorrent Network\par}
  \vspace{0.9em}
  {\normalsize The SwartzNet Project\par}
  \vspace{0.25em}
  {\normalsize July 2026\par}
\end{center}
\vspace{0.4em}

% ---- abstract -------------------------------------------------------------
\begin{center}
\begin{minipage}{0.86\textwidth}
\small\setlength{\parindent}{0pt}
\textbf{Abstract.} BitTorrent moves more data than almost any other
peer-to-peer system, yet it has no native way to \emph{find} anything:
discovery happens on centralized web indexes that are out-of-band,
fragile, and a single point of censorship. A purely peer-to-peer search
layer would let content be located directly by the same swarm that
serves it, with no coordinating server. The central problem is that a
search overlay normally demands its own protocol --- a new handshake
bit, a new DHT query, a new port --- which fragments the network and
makes the overlay trivial to fingerprint and block. SwartzNet solves
this by carrying search entirely inside protocol surfaces that mainline
clients already speak: the extension protocol, the DHT's signed-item
store, and ordinary torrents. A vanilla client sees nothing but standard
traffic. Search indexes are signed by their publishers, reconciled
between peers with a rateless set protocol, and defended against spam by
reputation and admission control rather than by a central authority. The
result is a distributed full-text index that is indistinguishable, on
the wire, from the network it lives in.
\end{minipage}
\end{center}
\vspace{1.1em}

\hypertarget{introduction}{%
\section{1. Introduction}\label{introduction}}

Finding a torrent means leaving the network that hosts it. A user
searches a web index, copies a magnet link, and only then rejoins the
swarm. Those indexes are centralized: they are operated by identifiable
parties, they are the point where takedowns and surveillance are
applied, and when one disappears its catalogue disappears with it. The
swarm itself --- millions of nodes that already store and exchange the
content --- holds no searchable record of what it carries.

What is needed is a search layer that is part of the network rather than
beside it: any node can publish what it indexes, any node can query its
peers and the DHT, and no node is required for the system to function.
The difficulty is doing this without changing what the network
\emph{is}. BitTorrent's reach comes from a shared set of standards (the
BEP series); a search system that introduces a new reserved bit or a new
DHT verb creates a second, incompatible network and paints a target on
every participant. SwartzNet's design goal is therefore not merely
``peer-to-peer search'' but \emph{invisible} peer-to-peer search:
full-text discovery that reuses existing standards so completely that a
node running it is, to any observer, a normal BitTorrent client that
happens to be unusually well-connected.

\hypertarget{design-constraint-mainline-compatibility}{%
\section{2. Design Constraint: Mainline
Compatibility}\label{design-constraint-mainline-compatibility}}

Every mechanism in SwartzNet is bound by one rule: a vanilla client must
see only standard traffic. Concretely, SwartzNet adds \textbf{no new
reserved handshake bit, no new DHT query, and no new UDP port.} It
builds only on established extension points:

\begin{itemize}
\tightlist
\item
  \textbf{BEP 3/5/9/10} --- the core protocol, the mainline DHT,
  metadata exchange, and the extension protocol (LTEP), which lets peers
  negotiate named sub-protocols inside the normal handshake.
\item
  \textbf{BEP 44} --- the DHT's store for small signed (mutable) and
  content-addressed (immutable) items, the substrate for a keyword
  index.
\item
  \textbf{BEP 46} --- mutable-torrent pointers in the DHT, reused as a
  stable pointer to a publisher's latest index snapshot.
\item
  \textbf{BEP 51} --- DHT infohash sampling, reused to discover what the
  network holds.
\end{itemize}

Because SwartzNet's peer-wire protocol is negotiated through LTEP, a
peer that does not advertise it is never sent a SwartzNet message;
because its DHT records are ordinary BEP-44 items, any client can read
them and none is required to understand them. Compatibility is not a
feature bolted on afterward --- it is the frame every other decision is
made inside. Any change that would let a vanilla client observe
non-standard traffic is, by definition, a defect.

\hypertarget{identity-and-signed-records}{%
\section{3. Identity and Signed
Records}\label{identity-and-signed-records}}

Search is only useful if results can be attributed and trusted, but a
network with no accounts cannot rely on identity being \emph{assigned}.
Each node instead generates a persistent ed25519 keypair the first time
it runs. That key is the node's publisher identity: it signs the records
the node contributes to the shared index, and its reputation --- good or
bad --- accrues to the key, not to an address.

A record is the atomic unit of the index: a signed statement that a
given keyword applies to a given infohash, at a given time, optionally
carrying a small proof-of-work stamp. Signing is designed so that
\textbf{the infohash is never disturbed.} When a publisher signs a
\texttt{.torrent}, the signature lives in top-level fields outside the
info-dictionary, so the content's identity is unchanged and a vanilla
client downloads it exactly as before. The signature authenticates
\emph{who is making a claim about} a torrent, never the torrent itself.
This separation --- stable content identity, attributable claims about
it --- is what lets an untrusted, open network accumulate an index whose
entries can still be weighed by their source.

\hypertarget{layer-l-the-local-index}{%
\section{4. Layer L: The Local Index}\label{layer-l-the-local-index}}

The foundation is entirely local. Every node maintains a full-text index
of the torrents it holds, built with an embedded search engine over the
torrent names and --- where the node chooses to index content --- the
text extracted from the files themselves (PDF, EPUB, office documents,
plain text, subtitles). This layer answers queries instantly and
offline; it is the ground truth a node can vouch for because it indexed
the bytes directly. Content extraction is opt-in per torrent, and
publication to the wider network can be turned off entirely, so the
operator controls both what text a node indexes and whether the node
announces anything beyond itself. The local index is deliberately
isolated: it knows nothing of the wire protocols above it, and the
layers above it never reach into its internals. Search results from
different layers are kept separate and reconciled only at the
presentation edge, never merged into a single ranked list that would
blur which layer --- and which trust model --- produced a hit.

\hypertarget{layer-s-peer-wire-search}{%
\section{5. Layer S: Peer-Wire Search}\label{layer-s-peer-wire-search}}

The first way a node's index reaches others is directly, over the peer
connections it already has. SwartzNet registers a named extension,
\texttt{sn\_search}, through LTEP. Immediately after the extension
handshake, a participating peer sends a capability announcement --- a
mask of which search services it offers.

The opt-in is enforced on the \emph{send} side, and structurally. When a
node observes a peer announce \texttt{sn\_search}, it mints an
unforgeable capability token for that peer; sending any SwartzNet frame
to a peer requires such a token, and there is no way to construct one
except from an observed announcement. A peer that announces nothing is
therefore never sent a frame --- the silence is a property of the code
paths, not a runtime check that could be bypassed or a policy that could
be misconfigured. Queries themselves carry a search term and a scope,
not a token; a responder answers according to its own sharing setting, a
per-peer rate limit, and a ban list, and is free to answer a peer it has
not itself queried. Results carry infohashes with names, sizes, and
seeder counts. Because the whole exchange rides inside an
already-established peer connection, it needs no new socket and is
invisible to peers who did not opt in.

\hypertarget{layer-d-the-dht-keyword-index}{%
\section{6. Layer D: The DHT Keyword
Index}\label{layer-d-the-dht-keyword-index}}

Peer-wire search only reaches nodes one is already connected to. To make
an index entry discoverable network-wide, SwartzNet publishes it into
the DHT as a BEP-44 mutable item. A keyword is mapped to a DHT address
by salting the publisher's key with the keyword, so each publisher owns
an independent namespace for every term and no two publishers collide. A
lookup for a term fetches, from each known publisher, the signed list of
infohashes that publisher associates with it, then merges and scores
those lists by publisher reputation.

BEP-44 items are small --- on the order of a kilobyte --- and expire
after a couple of hours. SwartzNet works within both limits: a publisher
re-announces on a timer, and when a keyword's entry would overflow the
size cap the oldest hit is evicted rather than the record being dropped
or a non-standard multi-item scheme being introduced. The keyword index
therefore rides the standard signed-item store exactly as any other
BEP-44 application would, and a vanilla client can read any of it. What
SwartzNet adds is not a new DHT behaviour but a \emph{convention} for
what the signed bytes mean.

\hypertarget{set-reconciliation}{%
\section{7. Set Reconciliation}\label{set-reconciliation}}

Two nodes that have each accumulated many signed records will overlap
heavily but not exactly. Naively exchanging full record sets to find the
difference wastes bandwidth proportional to what both already share.
SwartzNet instead reconciles record sets using a rateless invertible
protocol (RIBLT): each side streams coded symbols computed over its
records, and the \emph{difference} between the two sets is decoded from
a number of symbols proportional to the size of that difference, not to
the size of the sets. Neither side needs to know the difference size in
advance --- it keeps sending symbols until the other side can peel out
the missing records. This lets peers converge on the union of their
indexes cheaply, turning every peer connection into an opportunity to
fill gaps in coverage. Reconciliation is itself carried as
\texttt{sn\_search} messages, so it inherits the same opt-in,
mainline-invisible properties as ordinary queries.

\hypertarget{the-aggregate-index}{%
\section{8. The Aggregate Index}\label{the-aggregate-index}}

The per-keyword DHT index answers ``who claims this term?'' one
publisher at a time. A richer query --- a prefix scan, or a single
authenticated snapshot of a publisher's whole catalogue --- needs a
compact, verifiable structure larger than a single DHT item can hold.
SwartzNet's Aggregate index packs a publisher's signed records into a
signed B-tree: a self-contained, deterministic file with a signed
trailer, where every record carries its own signature and the whole tree
carries a fingerprint over its contents. A reader can verify the
trailer's signature before trusting anything, verify each record
independently, and run authenticated prefix queries over the tree ---
every answer carrying its own signature --- while the format's page
structure leaves partial, on-demand reads open for a later client.

The tree is distributed as an ordinary torrent, and a small BEP-44
pointer --- a ``publisher-pointer'' item at a fixed,
publisher-independent salt --- advertises its infohash together with the
tree's fingerprint. The fixed salt means every publisher's pointer lives
at an address derived the same way, so the DHT cannot be scanned for
SwartzNet users by salt; the fingerprint in the pointer binds the
pointer to a specific tree, so a reader who fetches the tree can confirm
its bytes are exactly the ones the signed pointer vouched for. This is
the general shape of the whole system: a signed, content-addressed
artifact distributed as a normal torrent, located through a standard
signed pointer. The Aggregate index ships off by default --- the
per-keyword index is the shipping path --- but rides the same seams, so
enabling it changes a mode, not the protocol.

\hypertarget{spam-resistance}{%
\section{9. Spam Resistance}\label{spam-resistance}}

An open index with no gatekeeper invites poisoning: a hostile publisher
can claim any keyword for any infohash. SwartzNet does not try to
prevent claims; it makes them cheap to \emph{discount}. Four mechanisms
compose:

\begin{enumerate}
\def\labelenumi{\arabic{enumi}.}
\tightlist
\item
  \textbf{Reputation.} Each publisher key accrues a Bayesian-smoothed
  score from the outcomes of the hits it has returned --- confirmed
  good, or flagged bad --- so a key that produces useful results is
  weighted up and a key that produces junk is weighted down. Smoothing
  means a handful of flags cannot instantly bury a key, and a handful of
  fresh keys cannot instantly vouch for one.
\item
  \textbf{Deny-by-default admission.} A previously unknown publisher is
  not trusted merely because a few other unknown publishers endorse it
  --- the classic Sybil move. New keys must clear a bar that fresh,
  mutually-endorsing identities cannot, or be seeded by an
  operator-provided anchor list.
\item
  \textbf{A known-good filter.} Confirmed-good infohashes are recorded
  in a compact Bloom filter, so results the local user (or its trusted
  set) has already vouched for can be surfaced preferentially at
  negligible cost.
\item
  \textbf{Optional proof-of-work.} A record may carry a hashcash stamp,
  letting a publisher spend computation to raise the cost of
  mass-producing claims. It is optional and bounded, a dial rather than
  a wall.
\end{enumerate}

Crucially, none of these decisions are made by a model or a prompt;
admission, reputation updates, and eviction are deterministic code with
explicit thresholds and fail-closed defaults. The index tolerates bad
actors instead of trying to exclude them, which is the only stance that
survives in a permissionless network.

\hypertarget{capability-negotiation-and-privacy}{%
\section{10. Capability Negotiation and
Privacy}\label{capability-negotiation-and-privacy}}

A node should share exactly what its operator intends and no more.
SwartzNet splits its capability mask in two: an operator-controlled half
(whether to answer queries at all, and whether to share file-name versus
content hits) and a daemon-owned half (runtime facts such as whether the
node is currently publishing). The two never mix, so an operator's
``save my sharing preferences'' can never accidentally toggle a runtime
fact. A node also enforces its \emph{current} sharing setting on every
inbound query, so tightening what it is willing to answer takes effect
at once. Unknown or future capability bits are tolerated, never
rejected, so the mask can grow without a flag day. Because sharing is
off unless announced and the send path is token-gated, the
privacy-preserving default is also the mainline-compatible one: a node
that shares nothing is wire-indistinguishable from a vanilla client.

\hypertarget{content-discovery}{%
\section{11. Content Discovery}\label{content-discovery}}

Two further mechanisms populate the index without a central catalogue. A
node may publish a \textbf{companion content-index}: a compact,
compressed snapshot of what it has indexed, distributed as an ordinary
single-file torrent and advertised through a BEP-46 mutable pointer.
Another node can \emph{follow} a publisher, fetch its snapshots as they
update, and import its catalogue --- verifying authorship and
de-duplicating across snapshots --- so coverage spreads by subscription
rather than by a server. Separately, a node can \textbf{crawl} the DHT
using BEP-51 infohash sampling: a bounded, polite worker pool that
samples infohashes from the nodes it reaches and expands its frontier
from their neighbours, discovering what the network holds so it can be
fetched, indexed, and re-published. Both mechanisms are opt-in and ride
only standard verbs; neither downloads or indexes anything the operator
did not ask for.

\hypertarget{conclusion}{%
\section{12. Conclusion}\label{conclusion}}

We have described a full-text search layer for BitTorrent that requires
no coordinating server and no departure from the standards the network
already runs on. Nodes index what they hold locally, answer each other
over existing peer connections, publish signed keyword records into the
DHT's standard signed-item store, reconcile those records with a
rateless set protocol, and --- optionally --- distribute whole signed
catalogues as ordinary torrents located through standard pointers. Trust
is placed in cryptographic identity and earned reputation rather than in
any authority, and spam is discounted by deterministic admission and
scoring rather than excluded by a gatekeeper. The binding constraint
throughout is mainline compatibility: every message rides an existing
extension point, every record is a standard DHT item, and a node that
shares nothing is invisible. The network needed no new protocol to
become searchable --- only a convention, signed and carried inside the
one it already speaks.

\hypertarget{references}{%
\section{References}\label{references}}

\begin{enumerate}
\def\labelenumi{\arabic{enumi}.}
\tightlist
\item
  B. Cohen, \emph{The BitTorrent Protocol Specification} (BEP 3).
\item
  \emph{DHT Protocol} (BEP 5); \emph{Extension for Peers to Send
  Metadata Files} (BEP 9).
\item
  \emph{Extension Protocol} (BEP 10).
\item
  A. Norberg et al., \emph{Storing Arbitrary Data in the DHT} (BEP 44).
\item
  \emph{Updating Torrents via DHT Mutable Items} (BEP 46).
\item
  \emph{DHT Infohash Indexing} (BEP 51).
\item
  A. Back, \emph{Hashcash --- A Denial of Service Counter-Measure}
  (2002).
\item
  S. Nakamoto, \emph{Bitcoin: A Peer-to-Peer Electronic Cash System}
  (2008).
\item
  L. Yang, Y. Gilad, M. Alizadeh, \emph{Practical Rateless Set
  Reconciliation} (SIGCOMM 2024) --- the rateless invertible-Bloom
  set-difference protocol.
\item
  SwartzNet architecture and protocol drafts:
  \path{docs/05-integration-design.md},
  \path{docs/06-bep-sn_search-draft.md},
  \path{docs/07-bep-dht-keyword-index-draft.md}.
\end{enumerate}

\end{document}
